\section{Lean formalization plan and the object-language endpoint}
\label{sec:formalization}

The archive contains three Lean files.  They are deliberately separated by
logical role.

All three files are under the archive's \path{Lean/} directory.
\begin{description}[style=nextline,leftmargin=0pt,labelindent=0pt]
\item[\code{SphericalKernel.lean}]
Coordinate-free finite algebra: vectors, dot/cross products, Gram and Lagrange
identities, the core spherical numerators, polar involution, face-area closure,
and the tetrahedral norm expansion.  It is intended to use only Mathlib's ring
tactics.
\item[\code{SphericalTrig.lean}]
An abstract constructive trigonometric system, a record of spherical core data,
and theorem interfaces for cosine, sine, haversine, half-angle, Delambre,
Napier, special-triangle, differential, area, and tetrahedral statements.
\item[\code{IPRABridge.lean}]
The concrete representation boundary for ProveIt's \code{ConciseIPRA}: codecs,
rapid names, apartness, fixed PR symbols, checker formulas, and the final
\code{Prf PRAAx} targets.
\end{description}

No Lean executable is installed in the delivery environment, so these files
were not compiled.  They contain an algebraic prototype and precise theorem
interfaces; they do not constitute a completed object-level \(\ipra\) proof.
The remaining tasks are explicit below.

\subsection{Stage 1: finite vector algebra}

Use a three-coordinate structure over an arbitrary commutative ring:
\begin{lstlisting}[language={}]
structure V3 (R : Type*) where
  x : R; y : R; z : R

def dot   (u v : V3 R) : R := ...
def cross (u v : V3 R) : V3 R := ...
def det3  (u v w : V3 R) : R := dot u (cross v w)
\end{lstlisting}
Prove by \code{ring}:
\begin{enumerate}
\item Lagrange:
\[
 (u\crossp v)\dotp(u\crossp w)
 =(u\dotp u)(v\dotp w)-(u\dotp v)(u\dotp w).
\]
\item Gram:
\[
 \detm(u,v,w)^2=\detm((u_i\dotp u_j)_{i,j=1}^3).
\]
\item The tangent numerator identities and
\((x-yz)^2+\tau^2=(1-y^2)(1-z^2)\).
\item Polar involution:
\((w\crossp u)\crossp(u\crossp v)=\tau u\).
\item The coordinate face-area closure and the expansion of its squared norm.
\end{enumerate}
These theorems require neither an order nor a field.  They are the most stable
part of the formalization and should be compiled first.  Run
\code{#print axioms} on the final declarations; no custom axioms should occur.

\subsection{Stage 2: an algebraic spherical core record}

To avoid square roots in the main identity proofs, define a record over a
commutative ring:
\begin{lstlisting}[language={}]
structure SphericalCore (R : Type*) [CommRing R] where
  x y z sa sb sc tau : R
  cA sA cB sB cC sC : R
  sa_sq : sa^2 + x^2 = 1
  sb_sq : sb^2 + y^2 = 1
  sc_sq : sc^2 + z^2 = 1
  gram  : tau^2 = 1 + 2*x*y*z - x^2 - y^2 - z^2
  cA_num : sb*sc*cA = x - y*z
  cB_num : sc*sa*cB = y - z*x
  cC_num : sa*sb*cC = z - x*y
  sA_num : sb*sc*sA = tau
  sB_num : sc*sa*sB = tau
  sC_num : sa*sb*sC = tau
\end{lstlisting}
The side cosine laws are just \code{cA_num} rearranged.  The angle cosine laws,
unit-pair equations, sine cross-products, five-part rules and Cagnoli identities
are proved by \code{ring} or \code{linear_combination} from these fields.  This
record is also the semantic target of the exact Python checker.

A second record adds order information only where needed:
\begin{lstlisting}[language={}]
structure ProperCore (R : Type*) [LinearOrderedField R]
    extends SphericalCore R where
  sa_pos : 0 < sa; sb_pos : 0 < sb; sc_pos : 0 < sc
  tau_pos : 0 < tau
\end{lstlisting}
It supports guarded cancellation and positive-root uniqueness.  The finite
algebra should remain in the weaker record whenever possible.

\subsection{Stage 3: abstract constructive trigonometry}

Reuse or import the class \code{TrigSystem} from the ProveIt trigonometric
packet.  Its data are an additive angle group, a commutative coefficient ring,
functions \code{sin}, \code{cos}, a quarter turn, the addition laws,
Pythagorean identity and quarter-turn values.  It already supports parity,
subtraction, periods, double-angle, product-to-sum and square half-angle laws.

Define a named spherical triangle with six angle terms and a core realization:
\begin{lstlisting}[language={}]
structure NamedSphericalTriangle (Angle R : Type*)
    [AddCommGroup Angle] [LinearOrderedField R]
    [TrigSystem Angle R] where
  a b c A B C : Angle
  core : ProperCore R
  cos_a : cos a = core.x
  sin_a : sin a = core.sa
  -- cyclic side equations
  cos_A : cos A = core.cA
  sin_A : sin A = core.sA
  -- cyclic angle equations
  half_side_positive : ...
  half_angle_positive : ...
\end{lstlisting}
Then formalize the formula atlas in dependency order:
\begin{enumerate}
\item core cosine/sine laws and polar substitutions;
\item half-square identities, followed by positive-root forms;
\item cross-multiplied Delambre identities;
\item Napier and tangent laws under denominator hypotheses;
\item cotangent, five-part, Cagnoli, right and quadrantal rules;
\item differential identities in a separate module parameterized by a
constructive derivative interface.
\end{enumerate}
The recommended theorem statements expose cross forms first.  Quotient theorems
should accept explicit hypotheses such as \code{cos x != 0} only at the ambient
field level; the object-level version must instead accept an apartness witness.

\subsection{Stage 4: codecs and primitive-recursive operations}

Pin the implementation to the inspected ProveIt revision
\code{a3c053cc6243}.  Within
\code{LeanProofs.ConciseIPRA}, construct actual \code{PR} syntax trees for:
\begin{enumerate}
\item pairing, projections and fixed-length tuples;
\item signed integers, positive-denominator rationals and dyadics;
\item addition, multiplication, powers of two, absolute value and comparison;
\item bounded sums, products, maximum, scan and finite-list folds;
\item three-vector dot, cross and determinant operations;
\item the bounded square-root scan (2.6);
\item finite exponential and arctangent sums with their explicit cutoffs.
\end{enumerate}
Derive the defining equations with \code{PRA.proj}, \code{PRA.comp},
\code{PRA.precZero} and \code{PRA.precSucc}.  Invalid rational codes should
decode to a fixed harmless default so that every symbol remains total.

\subsection{Stage 5: rapid-name arithmetic}

Represent a fixed real input by particular lower- and upper-endpoint
\code{PR 1} terms, an integer magnitude bound, and derivations of the interval
regularity formulas.  Do not make a Lean function \code{Nat -> Rat} the object
being quantified over.  For each operation, compile a new particular PR term
and prove a finite rational error estimate at a free precision index.

The essential numerical lemmas are:
\begin{align*}
 |xy-x'y'|&\le |x||y-y'|+|y'||x-x'|,\\
 |x^{-1}-y^{-1}|&\le 2|x-y|/d^2\quad(|x|,|y|\ge d/2),\\
 |\sqrt x-\sqrt y|&\le\sqrt{|x-y|},\\
 \left|\sum_{j=N}^{M}\frac{z^j}{j!}\right|&\le2^{-p}
 \quad\text{for the prescribed }N(B,p).
\end{align*}
Every induction invariant is a Boolean combination of natural equalities and
inequalities, encoded by characteristic functions, so it is quantifier free.
Only after the fixed-precision theorem is obtained should universal
introduction quantify over the precision.

\subsection{Stage 6: an object-language polynomial certificate checker}

Encode sparse multivariate polynomials as finite lists of integer coefficients
and exponent vectors.  Implement normalization by insertion, merging like
monomials and deleting zero coefficients.  Prove inside \(\PRAAx\):
\begin{enumerate}
\item evaluation commutes with polynomial addition and multiplication;
\item normalization preserves evaluation;
\item a normalized zero polynomial evaluates to zero;
\item each stored reduction certificate from the core ideal is sound.
\end{enumerate}
Prove these by structural primitive recursion and quantifier-free induction on
list lengths.  The external Python and Lean ring checks are specifications and
regressions.  Acceptance requires the internal soundness derivation, not merely
a Boolean checker returning \code{true}.

The existing \code{EqKernel}, \code{EqArith} and normalization modules in
\code{ConciseIPRA} should be reused where their interfaces match; otherwise the
small specialized normalizer above is sufficient.

\subsection{Stage 7: geometry-to-core realization}

Given fixed coordinate names for \(u,v,w\):
\begin{enumerate}
\item compute dot/cross/determinant names;
\item prove the unit-vector and Gram identities from the polynomial checker;
\item compute positive side-sine names with bounded square root;
\item use supplied lower bounds to construct reciprocals;
\item compute \((c_A,s_A)\), etc., and prove the unit-pair equations;
\item construct the angle names once using the two-chart \(\atanTwo\) routine;
\item prove their sine/cosine realization and branch intervals.
\end{enumerate}
The chart choice is made at one fixed certified precision and retained.  A
fresh undecidable branch choice at every output precision would not define a
Cauchy name.

\subsection{Stage 8: area and finite polygons}

Represent a geodesic polygon by a finite list of oriented great-circle normals
and a triangulation certificate.  Implement the lower/upper dyadic cell counts
in the equal-area chart (10.1).  The strip estimate (10.3) supplies the modulus.
Prove finite additivity at the cell level, then derive the exact lune formula.
The eight-region partition underlying Girard is finite combinatorics; encode
its incidence table directly.  Polygonal excess follows by primitive recursion
on the triangle list.

As a computationally lighter alternative for triangles, define the area name
by the solid-angle formula (10.12) and prove its agreement with the cell-area
name through Girard.  The solid-angle route is preferable for executable
geometry; the lune route establishes the intended surface-area semantics.

\subsection{Stage 9: tetrahedra}

The tetrahedral module needs only coordinate algebra and positive square
roots.  Construct the four oriented doubled face vectors, prove their sum is
zero by \code{ring}, define each dihedral cosine by a normalized negative dot
product, and expand the squared closure.  The sine-product identity imports the
already formalized planar sine law on three faces.  No spherical angle names or
area integration are required.

\subsection{Stage 10: the final derivability theorem}

For each fixed collection of input PR terms, assemble a closed formula
\(\mathsf{SphHyp}\to\mathsf{SphConclusion}\).  The final target has the shape
\begin{lstlisting}[language={}]
theorem spherical_atlas_derivable :
  Prf PRAAx [] (Fm.imp sphericalHypotheses sphericalConclusions)
\end{lstlisting}
where every auxiliary function symbol is a concrete compiled \code{PR} term.
After this theorem exists, apply \code{pra_theorem_iff} or
\code{pra_iff_native} to transfer it to the desired concise native
presentation.  The transport theorem is the final step, not a substitute for
the derivation.

\subsection{Validation commands}

In a checkout with Lean/Mathlib 4.32.0 and the pinned ProveIt revision, the
initial validation sequence is intended to be
\begin{lstlisting}[language=bash]
lake env lean Lean/SphericalKernel.lean
lake env lean Lean/SphericalTrig.lean
lake build ConciseIPRA.NativeBiInterpretation
lake env lean Lean/IPRABridge.lean
\end{lstlisting}
Resolve elaboration or tactic-version differences before treating the files as
checked.  Then audit the final ambient kernel theorems with \code{#print axioms}
and, separately, require an actual closed \code{Prf PRAAx} term for the
arithmetized theorem.  These are distinct acceptance criteria.
